% -----------------------------------------------------------------------------------
\section{Mixed policies and evidential semantics}\label{sec:mixed}

\subsection{A canonical mixed extension}

A pure causal matrix does not determine how prediction responds to a mixed policy, and continuity
alone does not specify the extension. One precise extension takes the predictor's response to be
linear in the policy and the action to be drawn independently of the response given the announced
distribution:
\[
  \mu_p=\sum_ip_i\mu_i,\qquad g(p)=p^\top M,\qquad V(p)=p^\top Mp ,
\]
where $g(p)_j$ is the causal value of option $j$ under the beliefs induced by $p$. This differs from
a predictor of the realised sample, which can instead yield the linear value $\sum_ip_iM_{ii}$. Define
\begin{description}[font=\normalfont,leftmargin=2.4em,labelwidth=1.6em,labelsep=.8em,itemsep=.35em]
  \item[$\ax{D}_m$] $p$ maximises $p^\top Mp$;
  \item[$\ax{J}_m$] every action in the support of $p$ maximises $p^\top M$;
  \item[$\ax{P}_m$] $p$ assigns probability zero to every weakly dominated column.
\end{description}

\begin{proposition}\label{prop:mixedPJ}
In the canonical extension, $\ax{P}_m$ and $\ax{J}_m$ are jointly satisfiable on every matrix.
\end{proposition}

\begin{proof}
For $0<\eta<1/|I|$ let $\Delta_\eta=\{p:p_i\ge\eta\text{ for all }i\}$. The correspondence
$p\mapsto\argmax_{q\in\Delta_\eta}p^\top Mq$ on the compact convex set $\Delta_\eta$ has nonempty convex
values and a closed graph, so Kakutani's theorem gives
$p^\eta\in\argmax_{q\in\Delta_\eta}(p^\eta)^\top Mq$. If column $a$ weakly dominates column $b$, full support
makes $(p^\eta)^\top Me_a>(p^\eta)^\top Me_b$, so $p^\eta_b=\eta$; otherwise moving mass from $b$ to $a$ would
improve the response. Let $p$ be a limit point as $\eta\downarrow0$. It assigns zero mass to every
dominated column. For any $q$, put $q^\eta=(1-|I|\eta)q+\eta\one\in\Delta_\eta$; then
$(p^\eta)^\top Mp^\eta\ge(p^\eta)^\top Mq^\eta$, and in the limit $p^\top Mp\ge p^\top Mq$. Hence $p$ is supported on
best responses to itself.
\end{proof}

\begin{proposition}\label{prop:mixedconflict}
$\{\ax{P}_m,\ax{D}_m\}$ and $\{\ax{J}_m,\ax{D}_m\}$ are inconsistent.
\end{proposition}

\begin{proof}
For $M=\bigl(\begin{smallmatrix}2&3\\0&1\end{smallmatrix}\bigr)$ and $p=(x,1-x)$, $V(p)=1+x$, so $\ax{D}_m$ selects
$x=1$, while column $2$ dominates and is the unique best response to every $p$, so $\ax{P}_m$ and
$\ax{J}_m$ select $x=0$.
\end{proof}

\begin{proposition}\label{prop:mixedN}
For rules that select a single lottery, $\ax{N}$ and $\ax{D}_m$ are inconsistent.
\end{proposition}

\begin{proof}
For $M=I_2$, the swap fixes $M$, so a neutral single lottery is $(\tfrac12,\tfrac12)$; but
$V(p)=p_1^2+p_2^2$ is maximised only at the two pure endpoints.
\end{proof}

A neutral \emph{correspondence} of maximising lotteries exists; a single neutral maximising lottery
need not. The pure consistency maps therefore cannot simply be relabelled as mixed maps.

\subsection{Contraction for lotteries}

For a rule $L$ that selects one lottery, the natural strengthening of contraction is
conditional-renormalisation consistency:
\[
  \ax{C}_m:\qquad L(M|_K)=\frac{L(M)|_K}{L(M)(K)}\quad\text{whenever }L(M)(K)>0 .
\]

\begin{proposition}\label{prop:mixedcirculant}
$\{\ax{N},\ax{R},\ax{C}_m,\ax{Q}\}$ is a minimal inconsistent family for single-lottery rules.
\end{proposition}

\begin{proof}
In the circulant extension of Theorem~\ref{thm:circulant}, neutrality forces the uniform lottery;
conditioning forces the uniform lottery on every embedded normalised matrix; row translations then
force the uniform lottery everywhere. For minimality, take the uniform distribution over the
choice set of, respectively, the responsive fixed-label rule, diagonal ranking, row-range ranking,
and universal indifference. The first two maximise a binary-independent weak order, so that
$G(M|_K)=G(M)\cap K$ whenever the right-hand side is nonempty, which is exactly what $\ax{C}_m$ requires of
uniform lotteries; row-range ranking is the witness that omits $\ax{C}_m$.
\end{proof}

Support contraction alone is much weaker: $p_i\propto1+r_i/r_{\max}$, with $r_i$ the row ranges and the
uniform lottery when $r_{\max}=0$, has full support everywhere, satisfies \ax{N}, \ax{F}, \ax{R}, and is
nonconstant. A complete mixed consistency map requires choosing among genuinely different output and
contraction semantics; the results above are the extensions established under the stated ones.

\subsection{Two evidential semantics}\label{sec:evsem}

The matrix does not supply an evidential criterion until the conditioning experiment is identified.
Two natural choices give different answers.

\emph{Reference-policy law.} With a full-support reference $\rho$ and $P(I=i,s)=\rho_i\mu_i(s)$,
conditioning on the policy index $i$ yields $\mu_i$, and the evidential value of $i$ is
$\E[U(s,i)\mid I=i]=M_{ii}$. This semantics gives diagonal optimality, and it is the matrix form of
\cri{6}$^\kappa$ on single-choice problems.

\emph{Conditioning on the realised action under a known mixed policy.} In the canonical extension the
environment has law $\mu_p$, and observing a positive-probability realised action does not change it,
so its conditional value is $g_i(p)$. With an explicit completion for zero-probability alternatives,
the resulting support-optimality condition is ratifiability.

The two condition on different variables under different joint laws, and neither can be silently
substituted for the other. This is one precise sense in which \enquote{evidential decision theory}
names a family rather than a criterion, as the literature on the tickle defence and on ratifiability
has long made clear \autocite{eells1982,jeffrey1983}.
